\begin{explanationbox}
	\textbf{\textcolor{CExplanation}{一句话直觉}}
	三次函数的零点个数可以通过极值点函数值的乘积符号快速判定。

	\medskip
	\textbf{\textcolor{CExplanation}{核心拆解}}
	\begin{description}[style=nextline, leftmargin=2.8em, labelsep=0.5em, itemsep=0.45em]
		\item[\textbf{要点一（极值点有无）：}]
			判别式\(\Delta'\)决定导函数是否变号：\(\Delta'\leq0\)时无极值点，函数单调；\(\Delta'>0\)时有两个极值点。
		\item[\textbf{要点二（乘积符号判零点）：}]
			极值点函数值乘积的符号直接影响零点个数：乘积负得三个零点，乘积零得两个零点（重根），乘积正得一个零点。
	\end{description}

	\medskip
	\textbf{\textcolor{CExplanation}{几何本质（极值点与高低点）}}
	极值点是图像上的山峰和山谷。若山和谷在\(x\)轴异侧，曲线三次穿过\(x\)轴；同侧则只穿一次；一个点在轴上则相切。

	\medskip
	\textbf{\textcolor{CExplanation}{代数意义（零点存在定理）}}
	由连续函数零点存在定理，函数在极值点区间端点的函数值异号时就有零点，乘积符号本质是区间端点符号的概括。

	\medskip
	\textbf{\textcolor{CExplanation}{考点价值（高考热点）}}
	\begin{itemize}[leftmargin=*, itemsep=0.5em, topsep=0.4em]
		\item \textbf{考法一（直接求零点个数）：} 已知具体函数，求导、算极值、乘积直接得零点个数。
		\item \textbf{考法二（含参零点分布）：} 已知零点个数或存在重根，反求参数范围，利用极值乘积等于零或与零的比较列式。
	\end{itemize}

	\medskip
	\textbf{\textcolor{CExplanation}{顿悟点}}
	\textbf{零点个数问题转化为极值点函数值的符号问题，无需解三次方程或画图，实现降维打击。}

	\medskip
	\textbf{\textcolor{CExplanation}{使用场景}}
	\begin{itemize}[leftmargin=*, itemsep=0.5em, topsep=0.4em]
		\item \textbf{场景一（函数零点判断）：} 遇到三次函数，先求导找极值点，再算乘积得零点个数。
		\item \textbf{场景二（参数讨论题）：} 含参三次函数零点个数为条件，利用极值乘积构造不等式或等式求解参数。
	\end{itemize}

\end{explanationbox}
